\begin{figure}[htbp]
\centering
\begin{tikzpicture}[x=1.05cm,y=1.05cm,>=Stealth,
  every node/.style={font=\small},
  coordinate curve/.style={draw=linkblue,line width=.7pt}]
  \fill[linkblue!5] (-3,0) circle (.72);
  \draw[coordinate curve] (-3,0) circle (.72);
  \node at (-3,0) {$\D$};
  \fill[linkblue!5] (0,0) ellipse (.86 and .65);
  \draw[coordinate curve] (0,0) ellipse (.86 and .65);
  \node at (0,0) {$\Omega$};
  \fill[linkblue!5] (2,-1.02) rectangle (4,1.02);
  \fill[white] (3,0) circle (.72);
  \draw[coordinate curve] (3,0) circle (.72);
  \node at (3,0) {$\D$};
  \node[font=\footnotesize] at (3,.88) {$|t|>1$};
  \draw[->] (-1.02,.22) -- node[above] {$f$}
    node[below,font=\footnotesize] {interior} (-2.08,.22);
  \draw[->] (2.02,.22) -- node[above] {$G$}
    node[below,font=\footnotesize] {exterior} (1.02,.22);
  \fill[linkblue] (-3,-.72) circle (1.5pt);
  \fill[linkblue] (0,-.65) circle (1.5pt);
  \fill[linkblue] (3,-.72) circle (1.5pt);
  \node[below] at (-3,-.76) {$\zeta$};
  \node[below] at (0,-.70) {$G(t)$};
  \node[below] at (3,-.76) {$t$};
  \draw[->,dashed] (3,-1.36) -- node[below] {$\zeta=f(G(t))$} (-3,-1.36);
\end{tikzpicture}
\caption{The two conformal coordinates, shown schematically; $\Omega$ need not
be an ellipse. The Cauchy projection uses the
exterior boundary parameter $t$ and its uniform circle measure. The
contraction representation uses the interior disk coordinate $\zeta$;
passing between them requires the positive Jacobian of the boundary map
$f\circ G$. The shaded exterior on the right is the domain of $G$.}
\label{fig:coordinates}
\end{figure}
